\documentclass[11pt]{article}
\usepackage[a4paper,margin=26mm]{geometry}
\usepackage{amsmath,amssymb,amsthm,booktabs}
\newcommand{\gen}[1]{\langle #1\rangle}
\newcommand{\Alt}{\operatorname{Alt}}
\newcommand{\Sym}{\operatorname{Sym}}
\setlength{\emergencystretch}{2em}
\newtheorem{theorem}{Theorem}
\newtheorem{lemma}[theorem]{Lemma}
\title{Finite end groups for alternating group constructions}
\author{Human Compression project}
\date{Working manuscript; distinct review pending}
\begin{document}\maketitle
\begin{abstract}
We recover four constructions of Fernandes and Leemans: rank six for
$n\equiv3\pmod4$, rank five in both even residue classes, and rank four
for $n\equiv0\pmod4$, with the source's degree restrictions and types.
We replace the degree-dependent intersection assertions by calculations
in four fixed permutation groups, using orbit partitions, specified
bijections between orbits, and parity. The passage from each finite
certificate to every degree is proved explicitly. The acquired rank
contraction theorem supplies the stated lower ranks as well. These are
components of the recovery of the source's principal existence theorem;
the remaining higher ranks and the degree-twelve rank-four case are
outside this manuscript.
\end{abstract}

\section{The permutations}
Write $n=m+10$, where $m\geq5$ and $m\equiv1\pmod4$. Permutations act on
$\Omega=\{1,\ldots,m+10\}$; unspecified points are fixed. Put
\begin{align*}
r_0&=(1\ 2)(3\ 4)\cdots(m-2\ m-1),\\
r_1&=(2\ 3)(4\ 5)\cdots(m-1\ m),\\
r_2&=(m\ m+1)(m+5\ m+8)(m+6\ m+9)(m+7\ m+10),\\
r_3&=(m+1\ m+2)(m+3\ m+5)(m+4\ m+6)(m+7\ m+10),\\
r_4&=(m+2\ m+3)(m+5\ m+8)(m+6\ m+7)(m+9\ m+10),\\
r_5&=(m+3\ m+4)(m+5\ m+6)(m+8\ m+9)(m+7\ m+10).
\end{align*}
These are the literal edges of the source's Theorem~4.13, with its
horizontal path numbered from the left and the six remaining vertices
numbered by rows. Each permutation is even: $r_0,r_1$ have $(m-1)/2$
transpositions, an even number, and the other four have four each.
They are distinct involutions. Generators at distance at least two
commute, as follows directly from the displayed disjoint transpositions.
The adjacent product orders are
\begin{equation}\label{type}
 |r_0r_1|=m,\quad |r_1r_2|=6,\quad |r_2r_3|=3,
 \quad |r_3r_4|=5,\quad |r_4r_5|=3.
\end{equation}
For the first equality, the product is an $m$-cycle on the initial path.
For the second, its restriction to $\{m-1,m,m+1\}$ is a three-cycle,
while its action on any earlier $r_1$-edge is a transposition. Such an
edge exists since $m\geq5$. The last three equalities involve just the
last eleven points and are verified either by direct multiplication or
by the finite certificate below.

\section{A generation lemma}
\begin{lemma}\label{generation}
Let $G=\gen{s_i:i\in I}$ be transitive on a finite set $\Omega$. Suppose
$G$ contains the group $\Alt(C)$ supported on a subset $C$ of at least
three points. If each $s_i$ fixes some point of $C$, then
$\Alt(\Omega)\leq G$.
\end{lemma}
\begin{proof}
First, if $A,B$ have size at least three and meet, then the two supported
groups $\Alt(A),\Alt(B)$ generate $\Alt(A\cup B)$. When $A\cap B$
contains at least two points this follows from the three-cycles based on those two
points. If the intersection has one point $a$, choose $b\in A\setminus B$
and a three-cycle in $\Alt(A)$ taking $a$ to $b$. Conjugating
$\Alt(B)$ gives $\Alt((B\setminus\{a\})\cup\{b\})$. The latter
set meets $B$ in at least two points; after merging those two sets,
the resulting set meets $A$ in at least $a,b$, reducing to the first case.
If one set contains the other the assertion is immediate.

Let $N$ be the normal closure of $\Alt(C)$ in $G$. The supports $C^g$
of its conjugates form a family of sets. The preceding observation shows
that each connected component of their intersection graph supports a
full alternating group in $N$. The corresponding unions, together with
any unused singleton points, are precisely the $N$-orbits. Let $D$ be
the $N$-orbit containing $C$. Normality makes $G$ permute the $N$-orbits.
Each generator fixes a point of $C\subseteq D$, so it preserves $D$.
Consequently $G$ preserves $D$, and transitivity gives $D=\Omega$.
Thus $N=\Alt(\Omega)$, as required.
\end{proof}

\section{The finite end and its faithful copies}
Set $K_j=\gen{r_1,\ldots,r_j}$, for $1\leq j\leq5$. Each orbit is contained in one of the following sets:
\[
 \{1\},\qquad P_a=\{2a,2a+1\}\quad
 (1\leq a\leq(m-3)/2),\qquad C=\{m-1,\ldots,m+10\}.
\]
On every $P_a$, only $r_1$ acts, and it interchanges the two points.
Every word therefore has exactly the same action on every $P_a$.
Keep $P_1$, the set $C$, and the fixed point $1$, and discard the
remaining pairs. Relabel the retained fifteen points by
\begin{equation}\label{identification}
 1\mapsto1,\quad 2\mapsto2,\quad3\mapsto3,\quad
 t\mapsto m+t-5\quad(4\leq t\leq15).
\end{equation}
Under this identification, the restricted generators $r_1,\ldots,r_5$
are exactly their $m=5$ versions. Restriction is faithful: a word fixing
the retained pair fixes every discarded pair as well. Thus this is an
isomorphism of marked groups, simultaneously for every $K_j$ and every
subgroup generated by a subset of these five marks. This proves a
uniform reduction to one finite end; no bound on $m$ is being inferred
from a list of examples.

The certificate uses the explicit permutations at $m=5$. It verifies
\begin{equation}\label{finiteend}
 |K_5|=12!,\qquad (K_5)_{2,3}=\Alt(\{4,\ldots,15\}),
\end{equation}
where the subscript denotes pointwise stabilization. It also checks
all $32^2=1024$ ordered pairs of subgroups generated by subsets of
$(r_1,\ldots,r_5)$ and obtains their intersection property. In particular
all the $K_j$ are C-groups for every $m$.

The graph of all six marks is connected, hence their group $G$ is
transitive. Each mark fixes a point of $C$; this is immediate from the
displayed permutations. By \eqref{finiteend} and the faithful
identification, $G$ contains the supported group $\Alt(C)$.
Lemma~\ref{generation} gives $\Alt(\Omega)\leq G$. All six generators
are even, so $G=\Alt(\Omega)$.

\section{Four intersection bounds}
Write $H_j=\gen{r_0,\ldots,r_j}$ and put $T=\{1,\ldots,m\}$.
The nontrivial orbit partitions needed below follow simply by following
the edges with labels at most $j$. All omitted points are singleton
orbits. They are as follows:
\begin{center}\small
\begin{tabular}{cl}
\toprule
Group & Nontrivial orbits\\\midrule
$H_1$ & $T$\\
$H_2$ & $T\cup\{m+1\}$;
 $\{m+5,m+8\}$, $\{m+6,m+9\}$, $\{m+7,m+10\}$\\
$H_3$ & $T\cup\{m+1,m+2\}$;
 $U=\{m+3,m+5,m+8\}$,\\
 & $V=\{m+4,m+6,m+9\}$; $\{m+7,m+10\}$\\
$H_4$ & $T\cup\{m+1,m+2,m+3,m+5,m+8\}$;\\
 & $Y=\{m+4,m+6,m+7,m+9,m+10\}$\\\bottomrule
\end{tabular}
\end{center}
There are two additional constraints, both checked directly on the
generators. Every element of $H_3$ commutes with
\[
 \tau=(m+3\ m+4)(m+5\ m+6)(m+8\ m+9).
\]
Equivalently, its actions on $U,V$ correspond under the three displayed
pairs. Every element of $H_4$ induces an even permutation on $Y$:
the restrictions of $r_2,r_3,r_4$ are double transpositions and those
of $r_0,r_1$ are the identity. Orbit preservation alone would not retain
these constraints.

For $j=2,3,4,5$, let $B_j\leq K_j$ preserve each orbit of $H_{j-1}$
setwise. For $j=4$ also impose commutation with $\tau$; for $j=5$
also impose even action on $Y$. Certainly
\begin{equation}\label{bound}
 K_{j-1}\ \leq\ H_{j-1}\cap K_j\ \leq\ B_j.
\end{equation}
These four groups $B_j$ have the same faithful fifteen-point models for
every $m$. Indeed, all discarded $P_a$ lie inside $T$, and every
$K_j$-element preserves each such pair. Discarding them therefore adds
or removes no orbit-preservation condition. The intersections of the
remaining orbits with the retained fifteen points are exactly their
$m=5$ counterparts under \eqref{identification}. The sets $U,V,Y$ and
the permutation $\tau$ are contained in $C$ and are carried to the same
fixed data as well.

The finite certificate computes the following exact subgroup equalities.
The middle column records what would remain if the additional two
constraints were omitted.
\begin{center}
\begin{tabular}{cccc}
\toprule
$j$ & Partition stabilizer order & Additional condition & $|B_j|=|K_{j-1}|$\\\midrule
2 &2& none&2\\
3 &12& none&12\\
4 &864& centralize $\tau$&144\\
5 &604800& even on $Y$&302400\\\bottomrule
\end{tabular}
\end{center}
The program tests equality as actual permutation subgroups, not merely
equality of their abstract isomorphism types. Consequently
$B_j=K_{j-1}$ uniformly, and \eqref{bound} proves
\[
 H_{j-1}\cap K_j=K_{j-1}\qquad(2\leq j\leq5).
\]
The pair $(r_0,r_1)$ is a rank-two C-group. By induction, the standard
end criterion for string C-groups now applies to $H_{j-1}$ and $K_j$.
It proves that each $H_j$, and in particular $H_5=G$, has the intersection
property. The criterion used here is McMullen--Schulte, Proposition~2E16;
it is an attributed common input.

\begin{theorem}
For every $n\equiv3\pmod4$, $n\geq15$, the displayed generators give a
rank-six string C-group representation of $A_n$ of type
$\{n-10,6,3,5,3\}$. The group $A_n$ also has representations of ranks
three, four and five.
\end{theorem}
\begin{proof}
Generation, the string relations, the type and the intersection property
have all been proved above. In \eqref{type}, the third, fourth and fifth
bonds have odd orders $3,5,3$. The acquired contraction theorem applied
successively gives ranks five, four and three of the same group. All
markings remain irreducible at each step, as required by that theorem.
\end{proof}

\section{Three families in even degree}
The same argument proves the remaining low-rank existence assertions
in even degree. We give the permutations, the faithful finite model,
and every additional intersection constraint. Let $e=(a\ b)$ be called
an edge of label $i$ when it is a disjoint factor of $r_i$. All generators
below are obtained by multiplying the edges of their label. The vertices
and labels specify the diagrams completely.

\subsection{Five marks in degree $n\equiv2\pmod4$}
For $n\geq10$, take
\[
 r_0=(1\ 2)(3\ 4),\quad r_1=(2\ 3)(4\ 5),\quad
 r_2=(3\ 4)(5\ 6),
\]
and give the path edge $(j\ j+1)$ label $3+(j\bmod2)$ for every
$6\leq j<n$. These are the source's Theorem~4.3 permutations. Their
type is $\{5,5,6,n-5\}$. The first two orders involve only five
points; the third has a three-cycle and at least one disjoint
transposition; the final product is an $(n-5)$-cycle on $\{6,\ldots,n\}$.
Every mark is even, and the string relations are immediate from the
edges.

For the end argument reverse the order of the marks, putting
$s_i=r_{4-i}$, and now use $H_j=\gen{s_0,\ldots,s_j}$ and
$K_j=\gen{s_1,\ldots,s_j}$. The finite end $K_4$ acts on
$C=\{1,\ldots,7\}$, the pairs
\[
 \{8,9\},\{10,11\},\ldots,\{n-2,n-1\},
\]
and the fixed point $n$. Only $s_1=r_3$ acts on these pairs, always by
the same interchange. Retain the first pair and the fixed point,
relabeling that fixed point as $10$. This gives a faithful marked
isomorphism to the $n=10$ end. The finite certificate proves its
intersection property by all $16^2=256$ ordered subset pairs and proves
that it contains the supported $\Alt(C)$. More precisely, $|K_4|=7!$
and its pointwise stabilizer outside $C$ is $\Alt(C)$.

The nonsingleton orbits of $H_1,H_2,H_3$ are respectively
\[
 \{6,\ldots,n\};\qquad
 \{3,4\},\ \{5,\ldots,n\};\qquad
 \{2,\ldots,n\}.
\]
For each $j=2,3,4$, the subgroup of $K_j$ preserving the preceding
partition is exactly $K_{j-1}$. In the ten-point model its order is,
respectively, $2,12,720$. Erasing the repeated pairs preserves all these
conditions, since they lie in the single long orbit in each partition.
The end criterion therefore proves the full intersection property for
every $n$. The full graph is connected and each mark fixes a point of
$C$; Lemma~\ref{generation} proves generation of $A_n$.
Reversing the type gives $\{n-5,6,5,5\}$. Its third and fourth bonds
are odd, so contraction gives ranks four and three as well.

\subsection{Five marks in degree $n\equiv0\pmod4$}
For $n\geq12$, put
\begin{align*}
 r_0&=(3\ 5)(4\ 6),&
 r_1&=(2\ 3)(6\ 7),\\
 r_2&=(1\ 2)(3\ 4)(5\ 6)(7\ 8),
\end{align*}
and give $(j\ j+1)$ label $3+(j\bmod2)$ for $8\leq j<n$.
These are the source's Theorem~4.5 permutations. They are even, satisfy
the string relations, and have type $\{3,4,6,n-7\}$. Again use
$s_i=r_{4-i}$. The fixed end has core $C=\{1,\ldots,9\}$, repeated
$s_1$-pairs $\{10,11\},\ldots,\{n-2,n-1\}$ and fixed point $n$.
Keeping the first pair gives the twelve-point model. Its end group has
order $9!$, has the intersection property by the $256$ subset-pair
checks, and contains the supported $\Alt(C)$.

The relevant orbit partitions are
\begin{align*}
H_1:&\quad \{8,\ldots,n\},\\
H_2:&\quad \{1,2\},\{3,4\},\{5,6\},\{7,\ldots,n\},\\
H_3:&\quad \{1,2,3,4\},\{5,\ldots,n\}.
\end{align*}
The action of $H_2$ on its first three pairs is the same interchange
on all three, or the identity on all three. In particular every element
commutes with both $(1\ 3)(2\ 4)$ and $(1\ 5)(2\ 6)$.
Every element of $H_3$ commutes with $(1\ 4)(2\ 3)$: the restrictions
of its generators to the first four points are the identity,
$(1\ 2)(3\ 4)$, and $(2\ 3)$. These constraints are identities on
actual permutations, not assumptions concerning abstract group types.

At stages $j=2,3,4$, the orbit-preserving subgroups of $K_j$ have
orders $2,24,2880$. Intersect the latter two with the indicated
centralizers. The resulting groups are exactly $K_{j-1}$, of orders
$2,12,960$. These equalities are checked in the twelve-point model.
The erased pairs lie in the long orbit and are disjoint from the three
specified centralizing permutations, so the equalities transport
faithfully to every degree. The end criterion proves the intersection
property. The supported $\Alt(C)$, connectivity and the fixed points
of each mark in $C$ again prove $G=A_n$ by Lemma~\ref{generation}.

\subsection{Four marks in degree $n\equiv0\pmod4$}
For $n\geq16$, number the vertices of source Theorem~4.4 by
\[
 (l,n,m,i,b,j,k,f)=(1,2,3,4,5,6,7,8),\qquad
 (g,a,d,h,c,e)=(9,10,11,12,13,14).
\]
The repeated letter $n$ on the left is the source's vertex name; it is
not a new degree parameter. The six remaining core edges and the path
are specified by
\begin{align*}
 r_0&=(1\ 2)(9\ 12)(10\ 13)(11\ 14),\\
 r_1&=(2\ 3)(4\ 12)(5\ 13)(11\ 14),\\
 r_2&=(3\ 4)(5\ 6)(7\ 8)(10\ 11)(13\ 14)
          \prod_{a=0}^{(n-16)/2}(15+2a\ \ 16+2a),\\
 r_3&=(4\ 5)(6\ 7)(9\ 10)(12\ 13)(11\ 14)(8\ 15)
          \prod_{a=0}^{(n-18)/2}(16+2a\ \ 17+2a).
\end{align*}
The last product is empty at $n=16$. The equations are an alternative
to the source's picture and extend its alternating path beyond vertex
$15$. Both tail counts have the parity needed to make all four marks
even. The string relations hold. The type is
$\{3,12,\operatorname{lcm}(n-8,6)\}$: the last two generators give
an alternating path on $n-8$ points and a separate component on six
points; the other two products are supported on the fixed core, with
possible repeated transpositions on the tail.

Set $s_i=r_{3-i}$. The finite end $K_3=\gen{r_2,r_1,r_0}$ acts on
$\{1,\ldots,14\}$ and on pairs
$\{15,16\},\{17,18\},\ldots,\{n-1,n\}$. On each pair only
$r_2$ acts, by the interchange. The restriction to the first sixteen
points is therefore faithful and gives the $n=16$ marked group.
It has order $144$ and the intersection property, verified on all
$8^2=64$ ordered subset pairs.

The nonsingleton $H_1$-orbits are
\[
 \{3,4,5,6,7,8\}\cup\{15,\ldots,n\},\qquad
 \{9,10,11,12,13,14\},
\]
with $1,2$ fixed. The group $H_2$ fixes $1$ and is transitive on its
complement, as its labelled graph shows. In the sixteen-point model,
the subgroups of $K_2,K_3$ preserving these respective partitions are
exactly $K_1,K_2$, of orders $2,24$. Discarding the extra pairs changes
neither condition. Thus the end criterion gives the intersection
property in every degree.

Generation has a different finite entry point in this family. Put
$u=r_1r_2r_3r_2$. On the core its nontrivial cycles of length greater
than two are
\[
 (2\ 6\ 3),\qquad(4\ 14\ 9\ 11\ 12),\qquad(5\ 10\ 13\ 8).
\]
On their complement $r_1$ fixes every point, and $u$ agrees with the
involution $r_2r_3r_2$. Hence, uniformly in $n$,
\[
 u^{12}=(4\ 9\ 12\ 14\ 11).
\]
Let $C=\{1,\ldots,6\}\cup\{9,\ldots,14\}$. This set is
$K_3$-invariant and its action there is independent of $n$. The finite
certificate takes the $144$ conjugates of the displayed five-cycle
under the sixteen-point end and proves that they generate precisely
the supported group $\Alt(C)$, of order $12!/2$. The same conjugates
are the same permutations of $C$ in every degree, so this is a uniform
inclusion $\Alt(C)\leq G$. Every mark fixes a point of $C$, and the
full graph is connected. Lemma~\ref{generation} again gives $G=A_n$.

This proves the rank-four assertion for all $n\equiv0\pmod4$,
$n\geq16$. The reversed type has third bond three, so one contraction
also gives rank three. The degree-twelve rank-four case is not covered
by this family and must be supplied separately in the principal recovery.

\section{The complete finite obligations of these three families}
The three end certificates above are finite mathematical premises,
verified by the full GAP program. For clarity their numerical checks
are collected here; their transfer to every degree was proved in each
subsection.
\begin{center}\small
\begin{tabular}{lrrrr}
\toprule
Family & Base degree & End order & Subset pairs & Final bounds\\\midrule
Five marks, $n\equiv2$ &10&5040&256&$2,12,720$\\
Five marks, $n\equiv0$ &12&362880&256&$2,12,960$\\
Four marks, $n\equiv0$ &16&144&64&$2,24$\\\bottomrule
\end{tabular}
\end{center}
The finite programs use actual pointwise and setwise stabilizers,
centralizers, and subgroup intersections. In the second row the
centralizers enforce the relations among the small orbits that would
be lost by an orbit-size calculation. In the last row the certificate
also verifies the pure five-cycle and its supported alternating normal
closure. No assertion of exhaustiveness of the constructions or their
Schl\"afli types is made.

\section{Attribution and the remaining obligation}
All four permutation families and their stated existence results belong to
Fernandes and
Leemans, Theorems~4.3--4.5 and 4.13 of the fixed edition arXiv:1810.12450v2,
4 January 2019. Their proof identifies several of the same end groups
and invokes computations whose extension to arbitrary degree is stated
without an explicit transport argument. Our contribution here is the
faithful deletion of repeated two-point actions, the finite upper bounds
with their diagonal and parity constraints, and the complete fixed
certificate. The construction itself is not presented as new. The
normal-closure argument uses the elementary supported permutation
method already used in our acquired common theory.

The rank contraction theorem is the acquired proof of
Brooksbank--Leemans, \emph{Rank reduction of string C-group
representations}; it is not reproved here. It changes the resulting
Schl\"afli types. Thus the lower-rank existence conclusions follow, while
no claim is made to recover the distinct exact diagrams or types of
source Theorems~4.11, 4.12 or 4.14. The present result is not by itself
the principal recovery of L038: the higher ranks and the degree-twelve rank-four case must still be supplied. No open-question, formalization,
measurement or priority claim is attached to this working component.

The uniform proof depends on the stated finite GAP calculation and the
standard end criterion. The separate runs at several degrees were
construction diagnostics only and are not premises of the argument.
The finite certificate checks four finite ends and their stated constraint subgroups;
its use does not assert that the infinite theorem is verified by a
bounded search in $n$.
\begin{thebibliography}{9}
\bibitem{FL} M. E. Fernandes and D. Leemans,
\emph{String C-group representations of alternating groups},
arXiv:1810.12450v2, 4 January 2019; Ars Math. Contemp. 17 (2019), 291--310.
\bibitem{BL} P. A. Brooksbank and D. Leemans,
\emph{Rank reduction of string C-group representations},
arXiv:1812.01055v2, 3 May 2019; Proc. Amer. Math. Soc. 147 (2019),
5421--5426. The acquired common proof and its distinct review are pinned
in the dependency manifest.
\bibitem{MS} P. McMullen and E. Schulte, \emph{Abstract Regular Polytopes},
Encyclopedia of Mathematics and its Applications 92, Cambridge University
Press, 2002, Proposition~2E16.
\end{thebibliography}
\end{document}
